\chapter{Réalisation de l'Application}

\section{Introduction}
Ce chapitre décrit la réalisation technique de \appname. Il présente l'environnement de développement,
puis l'architecture de la solution : le processus principal qui dialogue avec PostgreSQL, le pont de
communication qui expose les opérations à l'interface, et l'interface elle-même. Les sections suivantes
détaillent la sécurité, l'internationalisation, la génération des rapports et le type d'application
retenu.

\section{Environnement et Outils de Développement}
L'application a été développée dans l'environnement suivant :
\begin{itemize}
    \item Poste de développement : % À COMPLÉTER : modèle, processeur, mémoire, carte graphique
          poste de travail Windows x64 (caractéristiques matérielles à compléter) ;
    \item Système d'exploitation : Windows pour la compilation des paquets Windows, Linux x64 pour la
          compilation des paquets AppImage et Debian ;
    \item Environnement d'exécution : Node.js 22 ou supérieur pour la construction du client.
\end{itemize}
Les outils suivants ont été utilisés :
\begin{itemize}
    \item \textbf{Visual Studio Code} : éditeur de code, avec la configuration de projet présente dans
          le dossier \texttt{.vscode} ;
    \item \textbf{Electron 45} : moteur d'application de bureau, qui fournit la fenêtre, le moteur de
          rendu Chromium et l'accès au système de fichiers ;
    \item \textbf{Electron-Vite 5} et \textbf{Vite 7} : chaîne de construction des trois bundles
          (processus principal, pont de préchargement, interface) ;
    \item \textbf{React 19} et \textbf{React Router 7} : construction de l'interface et navigation ;
    \item \textbf{Tailwind CSS 4} : système de styles et jetons de conception ;
    \item \textbf{PostgreSQL 15 ou 16} : \ac{SGBD} relationnel qui héberge la base \texttt{eNote} ;
    \item \textbf{node-postgres 8} : pilote de connexion entre le processus principal et PostgreSQL ;
    \item \textbf{bcryptjs 3} : hachage des mots de passe ;
    \item \textbf{Electron Builder 26} : fabrication des installateurs Windows, AppImage, paquet Debian
          et archive Linux ;
    \item \textbf{lucide-react} et \textbf{Radix Dialog} : jeu d'icônes et composant de dialogue
          accessible.
\end{itemize}

\section{Architecture de l'Application}
L'architecture générale s'articule autour de l'infrastructure backend, de l'architecture frontend et de
spécifications techniques. Elle repose sur le modèle de processus d'Electron : une fenêtre ne parle
jamais directement à la base de données ; elle appelle une opération exposée par le pont de
préchargement, que le processus principal exécute après avoir vérifié les droits de l'appelant.

\subsection{Infrastructure Backend}
\subsubsection{Technologies Fondamentales}
Le processus principal est écrit en JavaScript moderne (modules \ac{ESM}, exécutés par Node.js intégré à
Electron). Il regroupe une quarantaine de services métier : authentification, personnes, profil de
l'élève, structure académique, emploi du temps, absences, observations, récompenses et sanctions,
examens, notes, documents, bibliothèque documentaire (catalogue, catégories, revue), messagerie,
notifications, rapports, statistiques, journal d'activité, administration des comptes et mise à jour du
client. Chaque service interroge PostgreSQL par des
requêtes paramétrées et renvoie des lignes normalisées en majuscules, ce qui donne aux services et à
l'interface un contrat stable.

La couche d'accès aux données gère un pool de connexions, ouvre une transaction pour les opérations qui
touchent plusieurs tables et relance la connexion enregistrée lorsque le serveur PostgreSQL redémarre.
Elle utilise également le mécanisme \texttt{LISTEN}/\texttt{NOTIFY} de PostgreSQL : un déclencheur écrit
une notification, le processus principal reçoit l'avis sur une connexion dédiée et pousse aussitôt
l'événement à la fenêtre, sans interrogation périodique.

\subsubsection{Architecture de la Base de Données}
La base s'appelle \texttt{eNote}. Elle est créée en encodage \ac{UTF-8} et appartient à un rôle
applicatif dédié, \texttt{enote\_app}, distinct du superutilisateur \texttt{postgres}. Le schéma
\texttt{public} contient 111 tables, une vue, 157 clés étrangères, 76 séquences, 34 déclencheurs et
51 fonctions. Les déclencheurs assurent l'horodatage des écritures, la production des notifications
et la tenue de compteurs dénormalisés ; les fonctions portent les règles
qui doivent s'appliquer quel que soit le chemin d'écriture.

Le déploiement de référence place le serveur à l'adresse \texttt{192.214.23.15}, sur le port
\texttt{5432}, et les fichiers partagés sur le partage \ac{SMB} \texttt{\textbackslash\textbackslash 192.214.23.15}\allowbreak\texttt{\textbackslash eNoteFiles}. La base ne stocke que des clés relatives pour les photos, les
documents et les captures d'anomalie : un poste de travail ne peut donc pas enregistrer un chemin local
qui n'existerait pas ailleurs. Les documents de cours sont classés par cursus
(\texttt{documents\textbackslash MAJOR\_CODE\textbackslash CHAPTER\_CODE\textbackslash
LESSON\_CODE}) et les documents de la bibliothèque par rayon
(\texttt{library\textbackslash TYPE\textbackslash CATEGORY\_CODE}).

Les trois tables de la bibliothèque documentaire (\path{library_category}, \path{library_item} et
\path{library_access}) sont créées à la première ouverture de la page, comme \path{document_access} :
une installation existante n'a donc pas de migration manuelle à exécuter. Le catalogue ne stocke que
des clés relatives, et les vidéos sont lues par un protocole interne \texttt{library-media://} qui ne
donné à la fenêtre aucun chemin du disque.\par

Les évolutions du schéma sont livrées sous forme de scripts numérotés dans
\texttt{deploy/server} (initialisation du rôle et de la base, création du compte administrateur,
contraintes, créneaux horaires, portée des secrétaires, droits de conduite, etc.), appliqués dans
l'ordre sur une base existante. Une sauvegarde \texttt{pg\_dump} est prise avant chaque campagne de
migration.

\subsubsection{Stockage Partagé des Fichiers}
Les photos, les documents de cours, les documents de la bibliothèque et les captures d'anomalie vivent
sur un partage \ac{SMB} du serveur, et non dans le dossier de chaque poste : une photo copiée sur le
poste A serait invisible au poste B et disparaîtrait à la réinstallation de A. Windows n'ouvre un
partage qu'à un compte qui possède déjà une session avec le serveur ; un poste jamais mappé lirait
donc le partage comme un dossier vide, et afficherait « fichier introuvable » alors que la base
répond normalement.

Le client ouvre cette session lui-même, avec le compte du préréglage de déploiement
(\texttt{smbUser}, \texttt{smbPassword}), au démarrage puis de nouveau, au plus toutes les vingt
secondes, lorsqu'un fichier reste introuvable. La session est établie par \texttt{net use} avec
l'option \texttt{/persistent:no} : aucun lecteur ni mappage ne subsiste dans le profil de
l'utilisateur, et un poste reconfiguré n'a rien à nettoyer. Les échecs sont nommés — identifiants
refusés, session concurrente d'un autre compte, partage absent, serveur injoignable, droits
insuffisants — puis consignés au démarrage, de sorte qu'une panne de droits ne se confonde pas avec
des fichiers perdus. Un poste Linux monte le partage par CIFS avant le lancement ; l'application
y voit alors un simple chemin local.

\subsubsection{Implémentation de la Sécurité}
La sécurité repose sur cinq mécanismes complémentaires.
\begin{itemize}
    \item \textbf{Hachage des mots de passe} : les mots de passe ne sont jamais stockés en clair ;
          bcryptjs les hache avec un coût de 10. La comparaison se fait par \texttt{bcrypt.compare}.
    \item \textbf{Requêtes préparées} : chaque requête passe par des paramètres positionnels
          (\texttt{\$1}, \texttt{\$2}, \ldots), ce qui neutralise l'injection \ac{SQL}.
    \item \textbf{Contrôle d'accès par rôle} : le rôle du compte est relu en base à chaque opération
          sensible. Le rôle scolaire d'une personne (qui elle est dans l'établissement) est conservé
          séparément et ne peut jamais tenir lieu d'autorisation. Un changement de rôle est une
          opération d'administration : il est refusé sur son propre compte et ne peut retirer le rôle
          Admin au dernier administrateur actif.
    \item \textbf{Règles temporelles paramétrables} : une option globale peut borner une écriture dans le
          temps, et le service concerné la relit en base à chaque appel. Le verrou de suppression des
          absences en est l'exemple : pendant les premières minutes d'une séance, un enseignant ne peut
          pas transformer une absence en présence, même en appelant le pont de préchargement
          directement.
    \item \textbf{Protection contre la devinette de mot de passe} : cinq échecs consécutifs suspendent
          le compte. Un administrateur le rétablit depuis la page des comptes suspendus, et la
          réinitialisation remet le compteur à zéro.
    \item \textbf{Journal d'activité} : les connexions, les refus, les changements de rôle et les
          modifications de données sont consignés dans \texttt{audit\_log}, en ajout seul. Les entrées
          ne sont ni modifiables ni supprimables depuis l'application, et celles de plus de douze mois
          sont élaguées automatiquement.
\end{itemize}
La fenêtre est créée avec l'isolation de contexte active et l'intégration Node désactivée ; seule
l'interface déclarée dans le pont de préchargement est exposée à la page. Les pièces jointes et les
photos passent par des dialogues natifs ouverts par le processus principal : l'interface ne lit jamais
un chemin arbitraire du disque. Les messages directs sont chiffrés au repos au moyen du coffre
\texttt{safeStorage} du système d'exploitation.

\subsection{Architecture Frontend}
L'interface est une application React 19 découpée en une centaine de composants. Elle s'appuie sur un
système de conception unifié : jetons de couleur, d'espacement, de rayon et de mouvement définis dans
\texttt{src/renderer/index.css}, composants partagés (boutons, dialogues, tableaux, pagination, états
vides, filtres, sélecteur de langue, bascule de thème) et une bibliothèque d'icônes unique.

La mise en page suit une coque stable : barre latérale fixe, en-tête de page, barre d'outils, zone de
contenu limitée à 1~440~px et barre d'enregistrement affichée uniquement lorsqu'un formulaire est
modifié. Les listes longues utilisent un tableau dense avec en-tête collant, sélection multiple,
actions groupées et pagination ; les dossiers utilisent un volet d'inspection à droite plutôt qu'une
page séparée. La densité, le thème et l'échelle de l'interface sont mémorisés par appareil.

Quatre apparences sont proposées (clair, sombre, noir absolu et suivi du système) et quatre palettes
d'accent (émeraude, rouge, or, bleu). Les polices Instrument Sans, Instrument Serif, JetBrains Mono et
Noto Naskh Arabic sont embarquées dans l'application, ce qui garantit un rendu identique sur un poste
sans Internet et dans les rapports imprimés.

\subsection{Internationalisation (i18n)}
L'application est disponible en anglais, en arabe et en français. Le choix de langue est une préférence
locale à l'appareil, conservée dans le stockage du navigateur, comme le thème : elle ne modifie pas le
compte et n'affecte pas les autres postes.

Le dispositif couvre trois niveaux. Pour l'interface, les dictionnaires arabe et français contiennent
3~790 entrées et l'anglais sert de langue de repli ; un changement de langue remonte l'arbre des
composants pour que chaque écran recharge ses données. Les dates et heures affichées suivent la même
langue : les noms de jours et de mois ne restent plus en anglais sur un écran arabe ou français. Pour la base, chaque libellé de référence porte
ses colonnes \texttt{\_en}, \texttt{\_ar} et \texttt{\_fr}, et chaque nom de personne ses colonnes
\texttt{\_la} et \texttt{\_ar} ; le processus principal substitue la variante active avant de renvoyer
les lignes. Lorsqu'une requête ne sélectionne que l'alias généré, un catalogue construit à partir de
\texttt{information\_schema} retrouve les variantes et les applique. Pour la direction d'écriture,
l'arabe bascule le document en \ac{RTL}, retourne la barre latérale et les tableaux, et les champs de
saisie multilingues affichent leurs variantes côte à côte.

\subsection{Génération et Intégration des PDF}
Les rapports ne dépendent d'aucune bibliothèque \ac{PDF} externe : ils sont produits par le moteur
Chromium d'Electron. Le processus principal construit un document \ac{HTML} autonome, l'ouvre dans une
fenêtre masquée, puis appelle \texttt{printToPDF} avec les marges et le format de page choisis. Les
polices embarquées sont déclarées dans la feuille de style du document, si bien que le fichier obtenu
est identique sur tout poste, y compris pour l'arabe.

Le rapport individuel de l'élève est le document principal : il reprend l'identité, la scolarité, les
affectations, les absences, les observations, les récompenses et les sanctions, dans la langue choisie
au moment de l'export. Le directeur sélectionne les rubriques à inclure, puis enregistre le fichier par
un dialogue natif. Les listes de séances et le journal d'activité peuvent également être exportés, le
premier au format imprimable, le second au format \ac{CSV}.

\subsection{Fonctionnement Hors Ligne}
Un serveur peut s'arrêter : coupure de courant, redémarrage de PostgreSQL, câble débranché. Le poste
continue alors de fonctionner sur ce qu'il a déjà lu, et le dit. Trois mécanismes complémentaires
couvrent la panne.

\begin{itemize}
    \item \textbf{Copies locales en lecture.} Les données lues par un poste — dossiers d'élèves,
séances, photos, emplois du temps, examens — restent dans son dossier de données avec la date de leur
dernière lecture. La fenêtre \emph{Enregistré sur cet ordinateur} indique ce qui est conservé, en
quelle quantité et depuis quand, et rappelle que la base redevient la référence dès qu'elle répond.
Les fichiers de la bibliothèque et des leçons peuvent en outre être enregistrés pour être ouverts
pendant une panne ; un cache limité à deux gigaoctets, élagué du plus ancien, évite de remplir le
disque.
    \item \textbf{File d'attente des écritures.} Les écritures additives — rapport d'anomalie, rapport
de test bêta, commentaire de notification — sont acceptées hors ligne et rejouées à la reconnexion,
dans l'ordre. Chaque type porte une clé naturelle qui rend le rejeu idempotent : une fiche de rapport
est un « upsert », un commentaire identique n'est pas enregistré deux fois. L'interface confirme
« enregistré sur cet ordinateur, sera envoyé au retour du serveur », et jamais « enregistré ».
    \item \textbf{Brouillons de formulaire.} Les formulaires longs — notes d'examen, fiche de l'élève,
signalement d'anomalie, détails d'une séance — conservent la saisie en cours, par formulaire et par
compte, pendant quatorze jours. À la réouverture, l'application propose de restaurer ou d'abandonner
le travail ; elle ne réapplique jamais un brouillon d'elle-même, afin qu'un poste partagé ne présente
pas la saisie d'un collègue. Le brouillon est effacé dès que l'enregistrement réussit.
\end{itemize}

Les écritures qui ne peuvent pas être rejouées sans risque restent refusées hors ligne : récompenses et
sanctions, affectations, planification, validation des notes et téléversement de pièces jointes. Le
message explique alors la règle au lieu d'échouer en silence. La connexion elle-même est bornée : seul
le dernier compte qui s'est connecté sur le poste peut se reconnecter sans serveur, avec son mot de
passe, pendant quatre-vingt-dix jours et cinq tentatives.

\subsection{Type d'Application (multi-pages / mono-page)}
\appname{} est une application de bureau à page unique (\ac{SPA}). La fenêtre charge un seul document
\ac{HTML} ; la navigation interne est assurée par React Router en mode \texttt{HashRouter}, ce qui
évite toute dépendance à un serveur \ac{HTTP} pour résoudre les routes. Ce choix découle du mode de
déploiement : il n'existe pas de serveur web dans l'architecture, et un rechargement complet de la page
n'a aucune raison de se produire. Le routeur conserve néanmoins des adresses distinctes pour chaque
écran, ce qui permet au fil d'Ariane, à la palette de commandes et aux liens de notification de désigner
une destination précise.

\section{Conclusion}
La réalisation repose sur trois couches clairement séparées : PostgreSQL pour la vérité des données,
le processus principal pour les règles et les contrôles, et une interface React pour la présentation.
Cette séparation explique que la sécurité ne dépende pas de l'écran affiché, que l'application puisse
fonctionner sans Internet, continuer à lire et à mettre en file ses écritures quand le serveur
s'arrête, et que les rapports soient générés avec les mêmes polices que l'interface. Le
chapitre suivant transforme ces choix en mode d'emploi, de l'installation du serveur à la journée de
travail de chaque rôle.
